\section{Related work and research positioning}
\subsection{Three complementary comparison perspectives}
We organize the evidence around three perspectives: theoretical formulas
and bounds on specified graph classes; heuristics and metaheuristics that
construct feasible solutions; and exact SAT/ILP methods that can establish
optimality upon completion. These categories are not mutually exclusive.
Exact solvers may use heuristics to obtain incumbents and theoretical bounds
to prune the search. At a time limit, an exact solver may return only bounds
or a feasible solution. Conversely, a heuristic solution attaining a proven
lower bound can also be certified optimal.

We assess SAT/ILP through formulation correctness, agreement with theory,
optimality coverage, remaining bound gaps, and computational cost. Without
a sufficiently systematic literature review, we do not claim that exact
methods are the least-studied class. The current experimental sequence is
to audit existing results and complete the SAT--ILP comparison before
evaluating an adapted genetic algorithm (GA).

\subsection{Labeling with two distance constraints}
Griggs and Yeh studied $L(2,1)$-labeling and graph-structural bounds
\cite{griggs}. Calamoneri's survey covers many $(h,k)$ regimes and graph
families \cite{calamoneri}. $L(2,1)$ is a classical case suitable as a default
and reference, not the most restrictive regime among all $L(h,k)$ problems.
On a fixed graph, increasing either threshold while holding the other fixed
can only reduce the set of feasible labelings.

For nontrivial trees, $\lambda_{2,1}(T)\in\{\Delta+1,\Delta+2\}$, and a
specialized linear-time algorithm exists \cite{hasunuma}. Fiala, Golovach,
and Kratochvíl proved NP-hardness of minimum-span labeling on trees when
$q$ does not divide $p$ \cite{fiala2008trees}. This motivates the choice
$(3,2)$ for investigation. The result concerns worst-case complexity over
an input class; it does not predict that every random tree or every subclass,
such as paths, is difficult. NP-hardness also does not imply that SAT and ILP
are the only exact methods. Excluding a general polynomial-time algorithm
requires $\mathrm{P}\ne\mathrm{NP}$ and the parameter domain of the theorem.

Murugan \cite{murugan2015corona} studied $L(2,1)$-labeling of corona products
and bounds for general coronas. Accordingly, this study does not claim to be
the first work on corona labeling. Substituting a family's maximum degree
into a general bound provides a theoretical reference, not automatically a new bound.

\subsection{Structural parameters and labeling variants}
Fiala et al.\ \cite{fiala2018parameterized} studied distance labeling and
channel assignment from a parameterized-complexity perspective. Their FPT
result for $L(p_1,\ldots,p_r)$ uses neighborhood diversity, $r$, and
$\max_i p_i$ jointly as parameters. We do not implement that FPT algorithm.
The result highlights the role of structure: runtime evaluation based solely
on vertex count is insufficient, and should also report degree information
and graph families or generation models.

Berthe et al.\ \cite{berthe2023edge} studied \emph{edge} labeling. This
corresponds to vertex labeling on a line graph, subject to the appropriate
distance convention. Their second condition applies to pairs joined by a
path of length 2, whereas the present implementation uses shortest-path
distance exactly 2. The definitions agree when $h\geq k$, but results cannot
be transferred unchanged to the API regime $h<k$. For example, $K_3$ with
$(h,k)=(1,2)$ admits labels $0,1,2$ under our convention, whereas the
common-neighbor condition requires three labels separated by at least 2.

King, Ras, and Zhou \cite{king2010trees} studied $L(h,1,1)$-labeling of trees,
requiring distinct labels at distances 2 \emph{and 3}. Their bounds cannot
therefore be substituted directly for $L(h,1)$ or $L(h,k)$ bounds here.
For example, labels $0,2,4,0$ on $P_4$ form a valid $L(2,1)$-labeling but
violate the distance-3 condition of $L(2,1,1)$. Edge labeling and distance-3
labeling are related extensions outside the present implementation scope.

\begin{table}[htbp]
\centering\small
\caption{Scope of related work and its role in this study.}
\label{tab:related}
\begin{tabular}{>{\raggedright\arraybackslash}p{0.24\linewidth}>{\raggedright\arraybackslash}p{0.31\linewidth}>{\raggedright\arraybackslash}p{0.36\linewidth}}
\toprule
Source & Problem & Role and limitations\\\midrule
Griggs--Yeh; Calamoneri \cite{griggs,calamoneri} & $L(2,1)$; $L(h,k)$ survey & Definitions, foundational results, and established references.\\
Fiala et al.\ \cite{fiala2008trees,fiala2018parameterized} & Tree complexity; structural parameters & Motivation for $(3,2)$ and structural diversity; not implemented solvers.\\
Berthe et al.\ \cite{berthe2023edge} & Edge / line-graph $L(p,q)$ & Related variant; uses the length-2-path convention.\\
King et al.\ \cite{king2010trees} & $L(h,1,1)$ on trees & Distance-3 constraints, outside the current model.\\
Vu Thanh et al.\ \cite{vu2026radio} & Radio labeling, preprint & Related formulation and SAT approach; not results of this study.\\
\bottomrule
\end{tabular}
\end{table}
Table~\ref{tab:related} distinguishes foundational references from extensions.
The shared term ``labeling'' does not make these problems or their optima identical.

\subsection{Relation to SAT methods for radio labeling}
The preprint by Vu Thanh et al.\ \cite{vu2026radio} uses order encoding and
incremental SAT for radio labeling; its evaluation sets the radio parameter
to the graph diameter. It is cited as a preprint, without assigning journal
or acceptance information in the absence of an official publication.
With radio parameter $r$, the threshold is $r+1-d(u,v)$. Setting $r=2$
yields $L(2,1)$, whereas $r=3$ also imposes a distance-3 constraint and thus
differs from $L(3,2)$. Spans and runtimes from the two studies are not compared
as results for an identical problem. Our approach uses two configurable
thresholds and a fresh solver at each span. This distinction specifies scope;
it does not itself establish novelty or superior performance.

\subsection{A genetic-algorithm heuristic baseline}
Sarhan et al.\ \cite{sarhan2026genetic} studied a GA for radio labeling with
the parameter set to the graph diameter. Their algorithm does not directly
solve $L(h,k)$. For example, $C_3$ has diameter 1 and radio span 2, whereas
$L(2,1)$ requires span 4. Their radio-labeling results are therefore not
imported into our comparison tables. One possible adaptation uses a GA to
search vertex permutations and a greedy decoder with thresholds $h$ at
distance 1, $k$ at distance 2, and 0 otherwise. Unconstrained vertices may
share labels. This supplementary baseline is a proposed design, not an
implemented or evaluated method in this manuscript, nor a reproduction
of the original algorithm and data.
